\properlane{the-promised-son}{How the Church waits for the Son already given}

\subsection{Whose Advent is Paul describing?}

There was a time when no one could recall the birth of Jesus. A promise about his coming in the flesh could then be held and handed on, but no parent could yet tell a child that he had been born. Thomas Aquinas asks his readers to take up the words of Romans 13 from that earlier standpoint. From the time when people first believed the promise there had been generations of believers. When the prophecies of his coming grew more frequent, someone who could look back over that long expectation might say with good reason that salvation was nearer than when they first believed. The meaning of the word \emph{we} has changed. The speaker represents the faithful who had hoped for Christ from the beginning. Thomas says expressly that it is the Church's use of the passage in Advent that allows him to hear Paul speak in their person.\footnote{Thomas Aquinas, \work{Super Epistolam B. Pauli ad Romanos lectura}, cap. 13, lectio 3, [86224], at \emph{Sed secundum quod Ecclesia haec verba adventus tempore legit} (on the Church's use of these words in Advent).}

Thomas had already explained what he took Paul to mean when he addressed the Romans. The salvation of eternal life is the end to which faith first directs a believer. Thomas considers two ways a person can come nearer to it. He can draw nearer in time to the end of his life. He can also grow better prepared to receive what he hopes for. Faith can issue in just action and the love with which a man acts can increase. Thomas illustrates readiness by the people admitted to the wedding when the bridegroom arrives. A length of time can be measured without knowing anything of the people who spend it. Readiness to enter when he comes requires a judgment about the people themselves. In Advent Thomas asks who else can say the words truthfully. He can give an answer about the approaching birth without making Paul entertain the thought that Jesus had yet to be born when he wrote to Rome.\footnote{Aquinas, \work{Super Romanos}, cap. 13, lect. 3, on Rom 13:11, including his appeal to Mt 25:10.}

This lets us hear a difference between two parts of the Introit. The psalmist speaks in the singular about what he has already done. He has trusted his God. Then he gives a reason in the plural, speaking of \emph{all} who wait. That word is broad enough to include people who expected the promised Son as well as those awaiting his return. The congregation that sings the Gradual can speak in their company. Sharing the expectation of an earlier believer does not require possessing more knowledge than that believer possessed about the terms in which God would fulfill it. One can ask to be taught God's ways precisely because one does not yet see how he will act. The Roman station named by the Missal for this Sunday is St Mary Major. To enter a church dedicated to Mary and sing those words gives the worshipper a particular woman to remember among the people who waited.\footnote{1962 Missal, \emph{Dominica prima Adventus}, p. [1], station heading and no. 1, and p. 2, no. 4.}

The Collect, by contrast, can name its addressee. The priest asks the Lord to come himself, and at the end of the prayer tells us of his eternal relation to the Father. Rescue is something one person does for another. The congregation calls upon someone they believe to be both able to perform it and willing to hear them. The psalm of the Alleluia lets Augustine develop that point. He asks what God's salvation is, and identifies it with his Christ. If a man asks to be given this salvation, there is a limit to how he can afterwards speak of his own part in obtaining it. Even the turning by which he has learned to want God is owed to the mercy for which he asks. On the lips of those who await Christ's birth, the two requests of the Alleluia are particularly exact. To show his people his mercy, the Lord can give them one whose face they may see.\footnote{Augustine, \work{Enarrationes}, on Ps 84, §6, at Ps 84:8 (85:7).}

\subsection{What grows of itself in a field}

The Communion antiphon can be taken as a plain promise to people who depend upon their land. God's favor will be evident in what their land can produce for them. Bellarmine reads the next-to-last pair of images in the psalm together. First truth has sprung from earth. Then the Lord will give goodness and the land will yield its fruit. The second saying explains the first for him. If that truth is the Son who has come from God, what does it mean for him to spring from earth? One kind of growth in a field comes of a person's ploughing and planting. Another can be seen in a wild flower. No husbandman put its seed into a furrow. It grows with what it receives from above, the sun and rain. Looking again at the two subjects of the Communion verse, Bellarmine understands the gift of the Lord as the sending of the Holy Spirit. He understands the earth that brings forth as the Virgin overshadowed by the Spirit. The fruit born of her is the Son.\footnote{Robert Bellarmine, \work{Commentary on the Book of Psalms}, Ps 84, paragraph 12 on the words of Ps 84:13, following v. 12; trans. John O'Sullivan (Dublin, 1866), p. 259.}

On this interpretation the fruit must truly be of the earth. The goodness sent from above does not displace the human mother who bears the child. The field Bellarmine pictures has had no human cultivator. Bellarmine can thus speak of one birth in which both things are true. Mary really bore the child, and she conceived him through the work of the Holy Spirit. Just to say that the fruit means a blessing from heaven loses what the psalm says about \emph{our earth}. To say simply that human beings gave the world a savior would leave the first verb in the chant unexplained.

Augustine also hears Christ born of Mary in the preceding verse. For him what Christ has taken from his mother is the mortal flesh in which he can offer himself. The immediate thing to marvel at is that truth has sprung from earth at all. Mortal human beings can lie, including about themselves. Yet the One who is Truth has taken what he could die in from a woman of that same earth. Augustine can go on to hear a second meaning in the image. When someone confesses his sin, he too is of the earth and he has begun to tell the truth. By the time Augustine reaches the Communion verse, his field has been opened up for the seed of God's word and needs the rain of grace. Bellarmine's untilled field and Augustine's furrows are both important to their owners. If one supplied the plough in Bellarmine's account in order to make the two fields look alike, one would take away his explanation of a virginal conception.\footnote{Augustine, \work{Enarrationes}, on Ps 84, §§10--12, at Ps 84:12--13 (85:11--12). Bellarmine, Ps 84, exposition 12.}

\subsection{The cloud in which the Son appears}

The coming foreseen in Luke is public. Men will see the Son of Man in a cloud with great power. To understand how this can belong beside the expectation of the Nativity, one needs to keep the difference between those occasions. At his birth he receives from his mother the humanity in which he can be seen. Jesus tells the disciples that in his coming they will see him in majesty. His birth makes the identity of the one who will come more wonderful, while the promise of his return tells them something more about the child.

Ambrose asks what sort of cloud one should imagine. His text speaks of clouds, in the plural, where the Missal has a single cloud. A rain cloud can make it difficult to see anything in daylight. In Scripture, however, God spoke to his people from a pillar of cloud. Ambrose considers the people through whom he made himself known. Moses, Joshua and the prophets can be called clouds because Christ speaks through them and the mystery becomes perceptible through human speech. The light cloud of the prophecy in Isaiah then suggests another way he comes. When Mary is told that the power of the Most High will overshadow her, the language of cloud and shade belongs to the conception of Jesus. Ambrose sees in the light cloud the Virgin. The prophets make Christ known in human speech. Mary bears the child of whom they spoke.\footnote{Ambrose, \work{Expositio evangelii secundum Lucam}, X.39--43, especially X.41--42, on the clouds, Is 19:1 and Lk 1:35. Compare Lk 21:27 in the Missal.}

Luke goes on to say more about the trees than one might remember from Matthew. In Matthew the disciples see the tender branch of the fig tree and know summer is near. Luke has the fig and \emph{all} the trees. Following the fruit in Luke, Ambrose imagines peoples who have come to confess God, including the Jewish people. A harvest of resurrection follows. But on looking at the fig tree from elsewhere in Scripture he also finds a reason to mistrust the apparent flourishing of leaves. The gardener in Jesus' other parable looks for fruit on a tree that has gone without it. Leaves can even be used to cover what a person is ashamed to show. So Ambrose allows a darker possibility. A man who has grown bold in doing wrong can give every appearance of success, and the approach of judgment will mean the end of that display. Ambrose turns to the other passages to make that distinction; in the words at Mass the trees are actually producing fruit.\footnote{Ambrose, \work{Expositio evangelii secundum Lucam} X.44--45, comparing Lk 21:29--30 with Mt 24:32 and Lk 13:6--9, and recalling Gen 3:7.}

It would be easy to make the land at Communion and all the trees in Luke mean the same thing on the strength of one common word. Bellarmine asks how the Son can be born of Mary. Ambrose asks what can be expected when faith begins to bear fruit among many peoples. The distinction lets a Christian hope for a real relation between the two. The one child born of the Virgin will be recognized by peoples from whom he could not have received his birth. When the Gospel has finished and the Church again addresses God in the voice of a single psalmist, each person offering the gifts can make an act of trust in the company of those other peoples. Guéranger observes something of that expectation at this very point of the Mass. As the bread and wine are offered, the Church continues to look toward him whom she expects.\footnote{Guéranger, \work{The Liturgical Year}, Advent, First Sunday of Advent, pp. 135--136, introduction to the Offertory.}

\subsection{The child whom the communicants await}

At the altar, the priest asks that the holy gifts may bring those who receive them with greater purity to their author. The faithful have already prayed to the Son to come. They can now wonder in what sense their approaching him is itself something his gifts can accomplish. What he needs if he is to approach with a pure heart is also something he can receive. The priest asks for precisely that cleansing through the holy gifts. At the Preface the priest summons the people to lift their hearts to the Lord and gives thanks to the Father. The Son whose coming the Advent congregation seeks shares the Father's majesty. The Holy Spirit who overshadowed Mary shares that majesty too. The charity with which a person answers the Lord's coming may have grown in him over the years, as Thomas says. That change in the worshipper did not give the Lord the power to come, either at his first coming or at this altar.

Guéranger opens his comment on the Communion chant by telling us when it is sung. The priest and people have received Communion. Yet the choir sings in the future tense about what the earth will give. In Guéranger's account the earth is Mary, made fruitful by the dew of heaven. The child she will bear has just given himself to the faithful in anticipation. There is no need to imagine a second birth at Christmas, or to instruct a communicant to forget having received the Lord in order to sing of his coming. What makes the words apt is the identity of the person. Someone who knows what Mary once received can recognize whom he receives, and can celebrate the approaching solemnity of that birth with living gratitude.\footnote{Guéranger, \work{Advent}, First Sunday of Advent, p. 137, the paragraph immediately before the Communion antiphon.}

The Postcommunion asks what remains for such a communicant to desire. Mercy is still something to receive. The Church asks to be able to greet the coming festivals with the honors due them. Thomas distinguished a person's progress through time from the action by which his charity grows. Here is one way that distinction can matter to people expecting Christmas. Paul has spoken of a debt of love to one's neighbor immediately before the lesson. There can be weeks still to pass before the feast, yet the time in which to go to the help of a particular neighbor is today. That work belongs to the person who now lives by faith in the Son already given. The day when he shall see the Son of Man with power remains a promise to be fulfilled.

\begin{fourSenses}
\item[Literal.] Israel asks God to show his favor, and hopes that the land will give a good yield. In Luke the fruit coming on the trees lets the disciples understand what is meant by the nearness of a season. Paul calls Christians who have come to faith to walk worthily now.
\item[Allegorical.] In Bellarmine's exposition the earth is Mary overshadowed by the Spirit, and the Son is the fruit she bears. Guéranger hears the same hope sung by Christians who have received that Son in Communion.
\item[Moral.] Thomas distinguishes the time elapsed from the growth of charity that makes a Christian ready to meet the Lord. The debt of love which Paul describes can require a person who awaits the feast to give his time to a neighbor in need of it now.
\item[Anagogical.] The peoples who have learned to confess the Son can hope for the harvest of resurrection of which Ambrose speaks. Those who receive Christ under the sacramental signs await the day they will see his majesty.
\end{fourSenses}
